\documentclass[11pt]{article}
\newcommand{\ListaTitulo}{Gabarito 01: Tipos de variável, tabelas e gráficos}
% Preâmbulo compartilhado: MATF14, Listas de Exercícios 2026
% Usado por ListaNN.tex e GabaritoNN.tex via % Preâmbulo compartilhado: MATF14, Listas de Exercícios 2026
% Usado por ListaNN.tex e GabaritoNN.tex via % Preâmbulo compartilhado: MATF14, Listas de Exercícios 2026
% Usado por ListaNN.tex e GabaritoNN.tex via \input{preamble}

\usepackage[utf8]{inputenc}
\usepackage[T1]{fontenc}
\usepackage[brazil]{babel}
\usepackage{fullpage}
\usepackage{amsmath,amssymb}
\usepackage{enumitem}
\usepackage{xcolor}
\usepackage{fancyhdr}
\usepackage{titlesec}
\usepackage{listings}
\usepackage{hyperref}
\usepackage{booktabs}
\usepackage{graphicx}
\usepackage{tikz}

\definecolor{matf14azul}{HTML}{0D3B54}
\definecolor{matf14dourado}{HTML}{C9971E}

\titleformat{\section}{\normalfont\Large\bfseries\color{matf14azul}}{\thesection}{1em}{}
\titleformat{\subsection}{\normalfont\large\bfseries\color{matf14azul}}{\thesubsection}{1em}{}

\providecommand{\ListaTitulo}{Lista de Exercícios}

\setlength{\headheight}{14pt}
\pagestyle{fancy}
\fancyhf{}
\lhead{\small MATF14: Estatística Econômica I}
\rhead{\small \ListaTitulo}
\cfoot{\thepage}
\renewcommand{\headrulewidth}{0.4pt}

\lstset{
  basicstyle=\ttfamily\small,
  breaklines=true,
  frame=single,
  columns=fullflexible,
  backgroundcolor=\color{gray!8},
  commentstyle=\color{matf14dourado},
  keywordstyle=\color{matf14azul}\bfseries,
  language=R
}

\newcounter{questaocounter}
\newenvironment{questao}[1]{%
  \stepcounter{questaocounter}%
  \bigskip\noindent\textbf{Questão #1.}\ %
}{\par}

\newcommand{\cabecalho}[2]{%
  \begin{center}
    {\Large\bfseries\color{matf14azul} #1}\\[0.3em]
    {\large #2}
  \end{center}
  \vspace{0.5em}
  \noindent\rule{\linewidth}{0.6pt}
  \vspace{0.5em}
}

\newenvironment{solucao}{%
  \par\smallskip\noindent\textit{\color{matf14dourado!80!black}Solução.}\ %
}{\hfill$\blacksquare$\par}


\usepackage[utf8]{inputenc}
\usepackage[T1]{fontenc}
\usepackage[brazil]{babel}
\usepackage{fullpage}
\usepackage{amsmath,amssymb}
\usepackage{enumitem}
\usepackage{xcolor}
\usepackage{fancyhdr}
\usepackage{titlesec}
\usepackage{listings}
\usepackage{hyperref}
\usepackage{booktabs}
\usepackage{graphicx}
\usepackage{tikz}

\definecolor{matf14azul}{HTML}{0D3B54}
\definecolor{matf14dourado}{HTML}{C9971E}

\titleformat{\section}{\normalfont\Large\bfseries\color{matf14azul}}{\thesection}{1em}{}
\titleformat{\subsection}{\normalfont\large\bfseries\color{matf14azul}}{\thesubsection}{1em}{}

\providecommand{\ListaTitulo}{Lista de Exercícios}

\setlength{\headheight}{14pt}
\pagestyle{fancy}
\fancyhf{}
\lhead{\small MATF14: Estatística Econômica I}
\rhead{\small \ListaTitulo}
\cfoot{\thepage}
\renewcommand{\headrulewidth}{0.4pt}

\lstset{
  basicstyle=\ttfamily\small,
  breaklines=true,
  frame=single,
  columns=fullflexible,
  backgroundcolor=\color{gray!8},
  commentstyle=\color{matf14dourado},
  keywordstyle=\color{matf14azul}\bfseries,
  language=R
}

\newcounter{questaocounter}
\newenvironment{questao}[1]{%
  \stepcounter{questaocounter}%
  \bigskip\noindent\textbf{Questão #1.}\ %
}{\par}

\newcommand{\cabecalho}[2]{%
  \begin{center}
    {\Large\bfseries\color{matf14azul} #1}\\[0.3em]
    {\large #2}
  \end{center}
  \vspace{0.5em}
  \noindent\rule{\linewidth}{0.6pt}
  \vspace{0.5em}
}

\newenvironment{solucao}{%
  \par\smallskip\noindent\textit{\color{matf14dourado!80!black}Solução.}\ %
}{\hfill$\blacksquare$\par}


\usepackage[utf8]{inputenc}
\usepackage[T1]{fontenc}
\usepackage[brazil]{babel}
\usepackage{fullpage}
\usepackage{amsmath,amssymb}
\usepackage{enumitem}
\usepackage{xcolor}
\usepackage{fancyhdr}
\usepackage{titlesec}
\usepackage{listings}
\usepackage{hyperref}
\usepackage{booktabs}
\usepackage{graphicx}
\usepackage{tikz}

\definecolor{matf14azul}{HTML}{0D3B54}
\definecolor{matf14dourado}{HTML}{C9971E}

\titleformat{\section}{\normalfont\Large\bfseries\color{matf14azul}}{\thesection}{1em}{}
\titleformat{\subsection}{\normalfont\large\bfseries\color{matf14azul}}{\thesubsection}{1em}{}

\providecommand{\ListaTitulo}{Lista de Exercícios}

\setlength{\headheight}{14pt}
\pagestyle{fancy}
\fancyhf{}
\lhead{\small MATF14: Estatística Econômica I}
\rhead{\small \ListaTitulo}
\cfoot{\thepage}
\renewcommand{\headrulewidth}{0.4pt}

\lstset{
  basicstyle=\ttfamily\small,
  breaklines=true,
  frame=single,
  columns=fullflexible,
  backgroundcolor=\color{gray!8},
  commentstyle=\color{matf14dourado},
  keywordstyle=\color{matf14azul}\bfseries,
  language=R
}

\newcounter{questaocounter}
\newenvironment{questao}[1]{%
  \stepcounter{questaocounter}%
  \bigskip\noindent\textbf{Questão #1.}\ %
}{\par}

\newcommand{\cabecalho}[2]{%
  \begin{center}
    {\Large\bfseries\color{matf14azul} #1}\\[0.3em]
    {\large #2}
  \end{center}
  \vspace{0.5em}
  \noindent\rule{\linewidth}{0.6pt}
  \vspace{0.5em}
}

\newenvironment{solucao}{%
  \par\smallskip\noindent\textit{\color{matf14dourado!80!black}Solução.}\ %
}{\hfill$\blacksquare$\par}


\begin{document}

\cabecalho{Gabarito 01}{Tipos de variável, tabelas de frequência e gráficos (Cap. 1, Aulas 02--03)}

\begin{questao}{1} \textbf{(Classificação de variáveis: pesquisa domiciliar)}
\begin{solucao}
\begin{enumerate}[label=(\alph*)]
  \item (i) número de pessoas: quantitativa discreta. (ii) faixa de renda: qualitativa ordinal
    (as faixas têm ordem natural, mas não distância numérica bem definida entre si). (iii) região:
    qualitativa nominal (sem ordem). (iv) renda exata em R\$: quantitativa contínua. (v)
    escolaridade: qualitativa ordinal.
  \item Ao agrupar a renda exata em faixas, perde-se a informação do valor preciso dentro de cada
    faixa (duas famílias na faixa "2 a 5 salários" podem ter rendas bem diferentes, mas ficam
    indistinguíveis). A vantagem prática é reduzir ruído de resposta (famílias relutam menos em
    indicar uma faixa do que um valor exato) e simplificar tabelas/gráficos de apresentação.
\end{enumerate}
\end{solucao}
\end{questao}

\begin{questao}{2} \textbf{(Tabela de frequências: setor de atividade)}
\begin{solucao}
Contando a lista: Comércio aparece 16 vezes, Serviços 16 vezes, Indústria 8 vezes ($n=40$).

\begin{enumerate}[label=(\alph*)]
  \item
  \begin{center}
  \begin{tabular}{lcc}
  \toprule
  Setor & $f$ & \% \\
  \midrule
  Comércio  & 16 & 40,0 \\
  Serviços  & 16 & 40,0 \\
  Indústria & 8  & 20,0 \\
  \midrule
  Total     & 40 & 100,0 \\
  \bottomrule
  \end{tabular}
  \end{center}
  \item Há empate entre Comércio e Serviços (ambos com 16 casos, 40\%), a amostra é
    \textbf{bimodal}.
  \item Um gráfico de barras (ordenado por frequência) é a escolha adequada: 3 categorias
    nominais, sem ordem natural, comparação de magnitude é o objetivo. Um gráfico de pizza
    também seria aceitável com apenas 3 fatias, mas é geralmente mais difícil de comparar
    precisamente duas fatias próximas (Comércio vs. Serviços) do que duas barras de mesma altura
    lado a lado.
\end{enumerate}
\end{solucao}
\end{questao}

\begin{questao}{3} \textbf{(Tabela de dupla entrada: mercado de trabalho)}
\begin{solucao}
\begin{enumerate}[label=(\alph*)]
  \item Fundamental/Médio: $70/250 = 28\%$ de desemprego. Superior: $40/250 = 16\%$ de desemprego.
  \item A manchete esconde que a taxa de desemprego é bem diferente entre os dois grupos de
    escolaridade (28\% vs. 16\%), a taxa agregada de 22\% é uma média ponderada que oculta essa
    heterogeneidade, relevante para qualquer política voltada a um grupo específico.
  \item Distribuição de escolaridade entre os desempregados: Fundamental/Médio $=70/110\approx
    63{,}6\%$; Superior $=40/110\approx 36{,}4\%$. Na população total, a distribuição é
    $250/500=50\%$ e $50\%$. As duas distribuições são diferentes: entre os desempregados, a
    proporção de baixa escolaridade é maior que na população geral, sinal de que escolaridade e
    situação ocupacional \textbf{não são independentes} nesta amostra (tema retomado no Capítulo 2).
\end{enumerate}
\end{solucao}
\end{questao}

\begin{questao}{4} \textbf{(Escolha de gráfico: série temporal vs. categorias)}
\begin{solucao}
\begin{enumerate}[label=(\alph*)]
  \item (i) Câmbio mês a mês: gráfico de \textbf{linha}, pois é uma série temporal e a ordem no tempo
    é informativa (mostra tendência, reversões, volatilidade). (ii) Empresas por região: gráfico
    de \textbf{barras}, pois região é uma variável qualitativa nominal sem ordem natural, barras
    permitem comparar magnitudes sem sugerir uma progressão inexistente.
  \item Ligar por linha 5 categorias sem ordem numérica (regiões em ordem alfabética) sugere
    visualmente uma "tendência" ou "trajetória" entre Centro-Oeste, Nordeste, Norte, Sudeste, Sul
    que não existe, a inclinação do segmento entre duas regiões consecutivas não tem
    significado algum, ao contrário da inclinação entre dois meses consecutivos em (i).
\end{enumerate}
\end{solucao}
\end{questao}

\begin{questao}{5} \textbf{(Dedução: frequências relativas somam 1)}
\begin{solucao}
\begin{enumerate}[label=(\alph*)]
  \item Por definição, $n = \sum_{i=1}^k f_i$ (soma das frequências absolutas de todas as
    categorias, que são exaustivas e mutuamente exclusivas, cada uma das $n$ unidades é contada
    em exatamente um $f_i$). Logo,
  $$
  \sum_{i=1}^k p_i = \sum_{i=1}^k \frac{f_i}{n} = \frac{1}{n}\sum_{i=1}^k f_i = \frac{n}{n} = 1.
  $$
  \item Suponha que a contagem seja feita em duplicidade: cada unidade é contada duas vezes, de
    modo que $f_i' = 2f_i$ para todo $i$, e portanto $n' = \sum_i f_i' = 2n$. A nova frequência
    relativa é
  $$
  p_i' = \frac{f_i'}{n'} = \frac{2f_i}{2n} = \frac{f_i}{n} = p_i,
  $$
  ou seja, $p_i'=p_i$: a frequência relativa é invariante a essa mudança de escala na contagem,
  porque o fator 2 aparece tanto no numerador quanto no denominador e se cancela. Por isso, \% é
  a medida correta para comparar grupos de tamanhos $n$ diferentes, contagens absolutas
  favoreceriam sempre o grupo maior, mesmo com composição interna idêntica.
\end{enumerate}
\end{solucao}
\end{questao}

\begin{questao}{6} \textbf{(Leitura crítica: gráfico enganoso)}
\begin{solucao}
\begin{enumerate}[label=(\alph*)]
  \item Num gráfico de barras, a informação visual transmitida é a \textbf{altura} (proporcional à área
    da barra, com largura fixa), o olho humano compara alturas relativas como se a base fosse
    zero. Se o eixo começa em R\$ 40.000 em vez de zero, uma diferença real pequena entre duas
    filiais (por exemplo, R\$ 42.000 vs. R\$ 45.000, uma diferença de $\approx 7\%$) ocupa quase
    toda a altura do gráfico, produzindo a impressão visual de uma diferença muito maior do que a
    proporção real entre os valores.
  \item Com o eixo começando em zero, as seis barras teriam alturas quase idênticas (todas
    próximas do topo do gráfico, já que a variação entre filiais é pequena frente ao nível geral
    de faturamento), a impressão visual correta seria de faturamentos parecidos entre as filiais,
    não de grandes disparidades.
\end{enumerate}
\end{solucao}
\end{questao}

\end{document}
